\documentclass[12pt, a4paper]{article}

% Pacchetti
\usepackage[utf8]{inputenc}
\usepackage[T1]{fontenc}
\usepackage[italian]{babel}
\usepackage{geometry}
\usepackage{titlesec}
\usepackage{enumitem}
\usepackage{fancyhdr}
\usepackage{xcolor}
\usepackage{booktabs}
\usepackage{array}
\usepackage{parskip}
\usepackage{hyperref}
\usepackage{longtable}

% Margini
\geometry{top=2.5cm, bottom=2.5cm, left=3cm, right=3cm}

% Colori
\definecolor{primary}{RGB}{30, 60, 120}
\definecolor{lightgray}{RGB}{245, 245, 245}
\definecolor{darkgray}{RGB}{80, 80, 80}

% Stile titoli sezioni
\titleformat{\section}
  {\large\bfseries\color{primary}}
  {\thesection.}{0.5em}{}[\titlerule]

\titleformat{\subsection}
  {\normalsize\bfseries\color{darkgray}}
  {\thesubsection.}{0.5em}{}

% Header e Footer
\pagestyle{fancy}
\fancyhf{}
\rhead{\small\textcolor{darkgray}{Verbale di Riunione}}
\lhead{\small\textcolor{darkgray}{MirrorFish Trading Offline}}
\rfoot{\small\textcolor{darkgray}{\thepage}}
\renewcommand{\headrulewidth}{0.4pt}

% Hyperlink
\hypersetup{
  colorlinks=true,
  linkcolor=primary,
  urlcolor=primary
}

% ─────────────────────────────────────────────
\begin{document}

% Intestazione / Titolo
\begin{center}
  {\LARGE\bfseries\color{primary} VERBALE DI RIUNIONE}\\[0.4cm]
  {\large Progetto: \textit{MirrorFish Trading Offline}}\\[0.2cm]
  \textcolor{darkgray}{\rule{0.6\textwidth}{0.4pt}}
\end{center}

\vspace{0.5cm}

% Tabella riepilogativa
\begin{center}
\begin{tabular}{>{\bfseries}p{4cm} p{9cm}}
\toprule
Data & 20 aprile 2026 \\
\midrule
Luogo & Discord (chiamata vocale) \\
\midrule
Partecipanti & Matteo, Marco, Davide \\
\midrule
Redatto da & Generato da trascrizione automatica \\
\bottomrule
\end{tabular}
\end{center}

\vspace{0.8cm}

% ─────────────────────────────────────────────
\section{Argomenti Discussi}

\subsection{Architettura del sistema e gestione dei componenti}

Il team ha discusso l'architettura generale del sistema MirrorFish, delineando una struttura modulare che distingue tra:

\begin{itemize}[leftmargin=1.5cm]
  \item \textbf{Componenti core}: il sistema principale responsabile della logica operativa fondamentale e delle regole non negoziabili.
  \item \textbf{Moduli di analisi e previsione}: estensioni che integrano capacità di Machine Learning per simulare scenari futuri sulla base di dati storici.
\end{itemize}

È emersa la necessità di bilanciare la robustezza delle fondamenta del sistema — basate su processi certi e deterministici — con la flessibilità delle componenti probabilistiche, che richiedono un approccio adattivo.

\subsection{Gestione delle decisioni e integrazione dati storici}

Il processo decisionale del sistema si basa sull'analisi di dati passati per simulare scenari futuri. Sono stati identificati i seguenti requisiti chiave:

\begin{itemize}[leftmargin=1.5cm]
  \item I dati storici devono essere \textbf{reperibili retroattivamente} a una data specifica, per garantire la coerenza tra addestramento e inferenza.
  \item La struttura dell'output deve essere \textbf{definita a priori e fissa}, in modo da essere compatibile con il contesto passato al modello in ogni momento.
  \item L'output finale non deve essere un semplice dato numerico, ma un \textbf{giudizio motivato} che integri la previsione con le regole operative stabilite.
\end{itemize}

\subsection{Incorporazione di capacità di Machine Learning}

È stato espresso un forte interesse nell'integrare modelli ML nel flusso decisionale. I punti di discussione principali sono stati:

\begin{itemize}[leftmargin=1.5cm]
  \item \textbf{Definizione dei confini funzionali}: occorre stabilire se la previsione ML debba essere un output separato e validabile oppure un fattore ponderato integrato nel flusso decisionale.
  \item \textbf{Ciclo di feedback}: è necessario definire un meccanismo che consenta al sistema di apprendere dai propri errori predittivi e auto-correggersi nel tempo.
  \item \textbf{Scalabilità}: va chiarito quali parti del sistema saranno \textit{hard-coded} (regole assolute) e quali saranno \textit{trainable} (regole variabili e aggiornabili).
\end{itemize}

\vspace{0.5cm}

% ─────────────────────────────────────────────
\section{To Do}

\begin{enumerate}[leftmargin=1.5cm, label=\textbf{\arabic*.}]

  \item \textbf{Analisti di Sistema — Documentazione dello Stato Base (MVP)}\\
  Formalizzare i processi non negoziabili e le regole operative fondamentali che costituiranno il primo rilascio funzionante del sistema. Il documento deve definire con precisione i confini del Minimum Viable Product.

  \item \textbf{Data Scientists — Specifiche dei Dati Predittivi}\\
  Definire esattamente quali dataset sono necessari per i modelli ML, specificando:
  \begin{itemize}
    \item Volume atteso dei dati;
    \item Formato di output richiesto dai modelli;
    \item Requisiti di retroattività storica.
  \end{itemize}

  \item \textbf{Architetti Software — Diagramma di Flusso Decisionale}\\
  Creare un diagramma UML o BPMN che mostri il percorso dei dati attraverso il sistema, evidenziando:
  \begin{itemize}
    \item I punti di decisione (\textit{Decision Points});
    \item I punti di intervento dell'intelligenza artificiale;
    \item Le interfacce tra componenti core e moduli ML.
  \end{itemize}

  \item \textbf{Team Tech Lead — Roadmap Tecnologica}\\
  Selezionare lo stack tecnologico appropriato per gestire sia la logica rigorosa (backend) sia l'elaborazione pesante dei modelli ML, tenendo conto di scalabilità e manutenibilità a lungo termine.

\end{enumerate}

\end{document}